% TOP_PROBLEM: {"id": 2831, "title": "Kirby Problem 3.33", "category_id": 7, "category": {"id": 7, "name": "topology", "display_name": "Topology", "description": "Properties preserved under continuous deformations.", "slug": "topology", "order_index": 7, "created_at": "2026-07-31T15:26:25.671Z"}, "status": "open", "problem_number": "KP-3.33", "set_id": 11}
% FIRST_POSTED: 2026-10-01
% ENTRY_KIND: statement_only
% UnsolvedMath ID 2831; source catalogue code KP-3.33
% !TeX program = lualatex
% Definition and sources only; this file is not an LLM solution attempt.
\documentclass[11pt,a4paper]{article}
\usepackage[margin=25mm]{geometry}
\usepackage{fontspec}
\setmainfont{lmroman10-regular.otf}[BoldFont=lmroman10-bold.otf,
 ItalicFont=lmroman10-italic.otf,BoldItalicFont=lmroman10-bolditalic.otf]
\usepackage{amsmath,amssymb,mathtools}
\usepackage{unicode-math}
\setmathfont{latinmodern-math.otf}
\AtBeginDocument{\let\setminus\smallsetminus}
\usepackage{xurl,hyperref}
\hypersetup{colorlinks=true,urlcolor=blue}
\setcounter{secnumdepth}{0}
\setlength{\parindent}{0pt}
\setlength{\parskip}{5pt}
\setlength{\emergencystretch}{3em}
\begin{document}
\section*{Kirby Problem 3.33}\label{2831}
{\small UnsolvedMath ID 2831 · KP-3.33 · Topology · Kirby's Problems in Low-Dimensional Topology}
\subsection{Definitions and mathematical statement}
For a closed connected orientable three-manifold $M$, a Heegaard splitting
is a decomposition $M=H_1\cup_\Sigma H_2$ into two handlebodies with common
boundary $\Sigma$. A genus-$g$ handlebody is a regular neighborhood of a
connected graph with first Betti number $g$, and its boundary is a closed
surface of genus $g$. The Heegaard genus $g(M)$ is the smallest genus of
such a splitting.

Choose orientations on $M$ and another closed connected three-manifold $N$.
A continuous map $f:M\to N$ has degree one when
\[
 f_*[M]=[N]\quad\text{in }H_3(N;\mathbb Z),
\]
where $[M]$ and $[N]$ are their fundamental classes. The conjecture is
\[
 \deg(f)=1\quad\Longrightarrow\quad g(M)\geq g(N).
\]
No assumption is made that $f$ is a covering, a bundle map, or injective on
fundamental groups. Although degree-one maps induce surjections on
fundamental groups, the number of group generators need not by itself
determine Heegaard genus. The task is to prove the displayed monotonicity
for all such maps or produce a degree-one map reversing it. The genera are
the minimal genera, rather than the genera of arbitrarily chosen splittings.
\subsection{Short English statement}
Can a degree-one map between closed oriented three-manifolds ever increase the minimum genus of a Heegaard splitting?
\subsection{Sources}
\begin{itemize}
\item[S1] ULAM.AI and the UnsolvedMath Contributors, supplied problems.json, record 2831 (KP-3.33), Kirby's Problems in Low-Dimensional Topology. Catalogue identity and source transcription; CC BY 4.0.
\url{https://huggingface.co/datasets/ulamai/UnsolvedMath}.
\item[S2] K3: A New Problem List in Low-Dimensional Topology, Problem 3.33. Expanded formulation of the supplied catalogue statement.
\url{https://math.berkeley.edu/sites/default/files/surv-295-ruberman-watermarked-author-pdf.pdf}.
\end{itemize}
\end{document}
